\section{Proper Scores and Pairwise Targets}\label{app:metrics}
The score-design method applies to a bounded known record-level target
$t(x,y)$, without requiring that it be a randomized classification gain.
For a label observed separately, write
\[
 t(x,y)=a(x)+yc(x),\qquad
 a(x)=\frac{t(x,1)+t(x,-1)}2,\quad
 c(x)=\frac{t(x,1)-t(x,-1)}2.
\]
On a finite domain, identification under arbitrary joint populations is
equivalent to $a,c\in\Hh_\A$. On the full binary cube with cap $b$, this is
$\deg(a),\deg(c)\le b$. Indeed the observation span on $(x,y)$ is
$\Hh_\A\otimes\operatorname{span}\{1,y\}$: a table may have separate
entries for the two labels. Its orthogonal complement gives indistinguishable
populations whenever the target lies outside that span. In particular the
common identification criterion for fixed forms and branching still applies.

For $z=(y+1)/2$ and $d=p_1-p_0$, the improvement in Brier loss is
\[
 t_{\rm Br}(x,y)=(p_0-z)^2-(p_1-z)^2
                =d(1-p_0-p_1)+yd.
\]
Thus degree at most $r$ for each probability function implies degree at most
$2r$ for this target's feature-only component and at most $r$ for its label
component. The bound can be strict because terms cancel. The actual degree,
or a certified approximation residual, determines the necessary cap.
Doubling is attainable: take $p_0=1/2$ and
$p_1=1/2+\gamma(x_1x_2+x_3x_4)$, with $0<\gamma\le1/4$.
The accuracy contrast has degree two, whereas the Brier feature-only
component is $-2\gamma^2(1+x_1x_2x_3x_4)$ and has degree four.
For a common depth-$b$ tree whose leaf values determine both probabilities,
the complete loss difference is still computable along that same path.

For log loss, clip each probability to a fixed interval
$[\eta,1-\eta]$, $0<\eta<1/2$, before defining the target. Then
\[
 \begin{split}
 a_{\log}&=\tfrac12\log\frac{p_1(1-p_1)}{p_0(1-p_0)},\\
 c_{\log}&=\tfrac12\log\frac{p_1(1-p_0)}{p_0(1-p_1)}.
 \end{split}
\]
These formulas estimate improvement of the clipped rules. A logarithm of
an ensemble average can introduce interactions of every order, even when
each probability has a low-degree representation. A common tree partition
again avoids that issue because the transformed values are constant on leaves.

The full-form program in Section~\ref{sec:pathdesign} replaces its
right-hand side by $\widehat t_y(S)$ for each label. Its pointwise identity,
population-specific variance objective, and invariance of unbiasedness to
calibration error remain unchanged. The symmetric-label minimax formula
$V(d,\A)$ and its sharp ratio were proved for $Yd(X)$; this extension of
the estimator does not assert that the identical formula is the global
minimax coefficient for every proper score. The path simulation does preserve
any bounded target's conditional mean and multiplies its conditional
second moment by at most the number of padded path sets.

\paragraph{AUC needs a pair design.}
Let $X^+\sim P(X\mid Y=1)$ and $X^-\sim P(X\mid Y=-1)$ be independent.
For a frozen score $s$, define
$k_s(x^+,x^-)=\mathbf1\{s(x^+)>s(x^-)\}
+\tfrac12\mathbf1\{s(x^+)=s(x^-)\}$.
The AUC difference is $\E[k_{s_1}-k_{s_0}]$ under the product of these two
conditional laws. Disjoint labeled case--control pairs give iid pair targets;
one can design unbiased scores on pair transcripts, retaining a separate
cap on each member. This differs from a single-record target: using all
case--control pairs introduces U-statistic dependence, and the sharp
single-record coefficient ratio cannot be asserted for that design.
Class prevalence, label-screening cost and the availability of both classes
must also enter an acquisition comparison. We establish the record-level
Brier/log-loss extension here and leave pair-design optimality unclaimed.

